\section{Exercise~15 and Tchoukaillon graphs}
\label{sec:ex15}

Exercise~15 assumes that the input has a perfect matching, $s=M=N$, and
looks for a graph on which $L$ increases by only~1 per round while each
round removes as few SAPs as possible. This section explains the
counting behind the exercise, where the Tchoukaillon numbers enter, and
why the definition of a Tchoukaillon graph can be read in more than one
way.

\subsection{Basic paths and parts (a) and (b)}

Fix a perfect matching $P$ and a current matching with $f$ free girls,
whose next dag has level~$L$. Recall that the basic paths are the $f$
vertex-disjoint augmenting paths in the symmetric difference of the
current matching with~$P$. We also need the number
\[
   \nu=\text{the largest number of vertex-disjoint SAPs of the current matching}.
\]
Always $\nu\le f$, since each SAP uses its own free girl.

\begin{lemma}[Exercise~15(a), (b)]
\label{lem:15ab}
\leavevmode
\begin{enumerate}[label=(\alph*)]
\item If every augmenting path of the current matching has level~$L$, then
  the dag satisfies
  \(f\le\min\{|G_0|,|B_1|,\ldots,|G_{L-1}|,|B_L|\}\).
\item Every maximal set of vertex-disjoint SAPs has at least
  $\ceil{\nu/(L+1)}$ members. In particular, if there are $f$
  vertex-disjoint SAPs (as there are under the hypothesis of~(a)), steps
  \Hs{11}--\Hs{17} find at least $\ceil{f/(L+1)}$ of them.
\end{enumerate}
\end{lemma}

\begin{proof}
(a) Under the hypothesis, the $f$ basic paths are SAPs. Let
$g_0\sg b_1\db g_1\sg\cdots\sg b_L$ be one of them. Its first part shows
that $g_j$ has distance at most~$2j$. If the distance were smaller, a
shortest directed path to $g_j$ followed by the rest of the SAP would be an
alternating walk from a free girl to a free boy with fewer than $2L-1$
edges; cutting out the closed loops of that walk would leave an
augmenting path of smaller level, which is impossible. So $g_j\in G_j$,
and likewise $b_j\in B_j$. Disjoint paths contribute different vertices
to each level.

(b) Let $S$ be a SAP. Its interior vertices are matched, and any
augmenting path through a matched vertex also contains that vertex's
mate. So a path that meets $S$ meets it at $g_0$, at $b_L$, or in one
of the $L-1$ couples of~$S$. Paths of a vertex-disjoint family can
share none of these, so $S$ meets at most $L+1$ members of the family.
Take a family of $\nu$ disjoint SAPs: each of them must meet some path of
the maximal set (otherwise it could be added), and the bound follows.
\end{proof}

The hypothesis of~(a) is stronger than what~(b) needs. This difference
between ``all augmenting paths have level~$L$'' and ``there are $f$
disjoint SAPs'' leads to the three readings of Section~\ref{sec:readings},
and the gap between $\nu$ and $f$ is what the worst cases of
Section~\ref{sec:staircase} exploit.

\subsection{The schedule and the Tchoukaillon numbers}
\label{sec:schedule}

A maximal matching has at least $\ceil{s/2}$ couples, because each of
the $s$ edges of $P$ must touch one of its couples and each couple
touches at most two of them. So~\Hs1 leaves at most $\floor{s/2}$ girls
free. Exercise~15(c) asks for a graph that starts with this largest
possible number of free girls and, at every level~$L=2,3,\ldots$,
removes only the number of SAPs that Lemma~\ref{lem:15ab}(b) forces when
$\nu=f$. The numbers of free girls and of SAPs removed then follow the
\emph{Tchoukaillon schedule}
\begin{equation}
  f_2=\floor{s/2},\qquad a_L=\ceil{f_L/(L+1)},\qquad
  f_{L+1}=f_L-a_L=\floor{Lf_L/(L+1)},
  \label{eq:schedule}
\end{equation}
which stops when $f$ reaches~0. For $s=16$ it reads:
\begin{center}
\begin{tabular}{c|ccccc}
\toprule
level~$L$ & 2&3&4&5&6\\
free girls~$f_L$ & 8&5&3&2&1\\
SAPs removed~$a_L$ & 3&2&1&1&1\\
\bottomrule
\end{tabular}
\end{center}
With~\Hs1 counted as round~1, that is six rounds.

Let $T(s)$ be the number of rounds of the schedule, that is, the last
level~$L$ with $f_L>0$ (and $T(1)=1$, $T(0)=0$). Answer~15(c) observes
that~\eqref{eq:schedule} is the recurrence of exercise~11(b) for the
solitaire game Tchoukaillon: in Knuth's notation there, $f_L=l_{L-1,s}$
and $a_L=m_{L,s}$, by parts (ii) and~(i). Consequently
\begin{equation}
  T(s)=n\quad\iff\quad \Ch_n\le s<\Ch_{n+1},
  \label{eq:Tdef}
\end{equation}
where the \emph{Tchoukaillon number} $\Ch_n$ is the least number of
stones for which the winning strategy uses pit~$n$:
\[
\begin{array}{c|cccccccccccccc}
 n&1&2&3&4&5&6&7&8&9&10&11&12&13&14\\\hline
 \Ch_n&1&2&4&6&10&12&18&22&30&34&42&48&58&60
\end{array}
\]
One can also see~\eqref{eq:Tdef} directly. Run the schedule backwards:
the last round has $f_n=1$, and the smallest $f_{L}$ that leads to a
given $f_{L+1}$ is $\ceil{(L+1)f_{L+1}/L}$. This gives Knuth's formula
$\Ch_n=\bigl\lceil2\bigl\lceil\frac32\bigl\lceil\frac43\cdots
\bigl\lceil\frac n{n-1}\bigr\rceil\cdots\bigr\rceil\bigr\rceil\bigr\rceil$
of exercise~11(c). Exercise~11(d) shows that
$\Ch_n=n^2/\pi+O(n^{4/3})$, so $T(s)=\sqrt{\pi s}+O(s^{1/6})$.

These are statements about the \emph{numbers} in~\eqref{eq:schedule}.
Whether any graph realizes them is the question of exercise~16.

\subsection{Knuth's sixteen-pair example}
\label{sec:knuth16}

Answer~15(d) describes a graph with $s=16$ by a sequence of
matrices (Fig.~\ref{fig:knuth16rows}). Label the edges of the perfect
matching $g_k\sg b_k$, $1\le k\le16$. A row such as $1\ 2\ 4$ stands for
the basic path $g_1\sg b_1\db g_2\sg b_2\db g_4\sg b_4$: its cells are the
pairs it visits, and moving one cell to the right means following a
couple. Each matrix lists the basic paths before one round. All its rows
have $L$ cells, so they are $f$ disjoint SAPs and $\nu=f$. The graph has
the 16 edges $g_kb_k$ and 19 more, listed in~\eqref{eq:knuth16edges} of
Appendix~\ref{app:knuth16}.

An edge $g_ib_j$ with $i\ne j$ jumps from pair~$i$ to pair~$j$. Every SAP
that Algorithm~H finds in this graph moves through the current matrix by
following couples to the right along a row and by such jumps from row to
row within a column. For
example, the first round finds the SAP $g_5\sg b_3\db g_4\sg b_2$: it
starts at the first cell of row~$5\ 6$, jumps to the first cell of
row~$3\ 4$, moves right along that row, and jumps up to the last cell of
row~$1\ 2$, whose boy is free. Augmenting replaces these three rows by
the two rows~$1\ 2\ 4$ and $3\ 5\ 6$ of the next matrix. In general a SAP
that meets $k$ rows leaves $k-1$ rows, each with one more cell.

\begin{figure}[ht]
\centering
% Basic-row matrices of answer 15(d).  Rows of one colour are merged by a
% single SAP of the next round (green: a row that a SAP consumes whole).
\begin{tikzpicture}[
   cellbox/.style={draw=black!25,minimum width=4.6mm,minimum height=4.3mm,
     inner sep=0pt,font=\scriptsize},
   hd/.style={font=\footnotesize,anchor=south},
   mstep/.style={-{Stealth},black!45},
   steplab/.style={font=\tiny,text=black!60,below}]
\def\w{.48}
\newcommand{\mrow}[4]{% x, y, list, colour
  \foreach \v [count=\i] in {#3}{
    \node[cellbox,fill=#4!35] at ({#1+(\i-1)*\w},{#2}) {\v};}}
\def\h{.43}
% L = 2
\node[hd] at ({0+.5*\w},.3) {$L=2$, $f=8$};
\foreach \a/\b/\c [count=\k from 0] in {1/2/rowA,3/4/rowA,5/6/rowA,7/8/rowB,9/10/rowB,11/12/rowB,13/14/rowC,15/16/rowC}
  {\mrow{0}{-\k*\h}{\a,\b}{\c}}
\draw[mstep] (1.0,-1.4)--(1.75,-1.4) node[steplab,midway] {3 SAPs};
% L = 3
\def\x{2.05}
\node[hd] at ({\x+\w},.3) {$L=3$, $f=5$};
\mrow{\x}{0}{1,2,4}{rowA}
\mrow{\x}{-\h}{3,5,6}{rowA}
\mrow{\x}{-2*\h}{7,8,10}{rowA}
\mrow{\x}{-3*\h}{9,11,12}{rowA}
\mrow{\x}{-4*\h}{13,14,16}{rowC}
\draw[mstep] (3.55,-.65)--(4.3,-.65) node[steplab,midway] {2 SAPs};
% L = 4
\def\x{4.6}
\node[hd] at ({\x+1.5*\w},.3) {$L=4$, $f=3$};
\mrow{\x}{0}{1,2,4,6}{rowA}
\mrow{\x}{-\h}{3,5,8,10}{rowA}
\mrow{\x}{-2*\h}{7,9,11,12}{rowA}
\draw[mstep] (6.6,-.43)--(7.35,-.43) node[steplab,midway] {1 SAP};
% L = 5
\def\x{7.65}
\node[hd] at ({\x+2*\w},.3) {$L=5$, $f=2$};
\mrow{\x}{0}{1,2,4,6,10}{rowA}
\mrow{\x}{-\h}{3,7,9,11,12}{rowA}
\draw[mstep] (10.1,-.2)--(10.85,-.2) node[steplab,midway] {1 SAP};
% L = 6
\def\x{11.15}
\node[hd] at ({\x+2.5*\w},.3) {$L=6$, $f=1$};
\mrow{\x}{0}{1,2,4,6,10,12}{rowA}
\end{tikzpicture}

\caption{The matrices of answer~15(d). Each row is a basic path, given by
the indices of the pairs of the perfect matching that it visits. Within
each matrix, rows of the same colour are merged by one SAP of the next
round (colours are assigned separately in each matrix); the green row at
$L=3$ is itself a SAP, and augmenting along it matches all three of its
pairs by $P$.}
\label{fig:knuth16rows}
\end{figure}

The first matrix says that the matching after~\Hs1 consists of the
couples $b_1g_2$, $b_3g_4$, \dots, $b_{15}g_{16}$. But the graph also
contains the edges $b_5g_8$ and $b_4g_6$, and so this matching has
augmenting paths of two different levels starting at the same free girl
(Fig.~\ref{fig:knuth16paths}):
\begin{equation}
  \pi_2=g_7\sg b_7\db g_8\sg b_8\quad(\text{level }2),\qquad
  \pi_3=g_7\sg b_7\db g_8\sg b_5\db g_6\sg b_4\quad(\text{level }3).
  \label{eq:PQ}
\end{equation}
So the hypothesis of~15(a), that every augmenting path has the same
level, already fails at the first round. On the other hand,
Appendix~\ref{app:knuth16} gives explicit input orders for this graph
under which Algorithm~H follows the schedule $8,5,3,2,1,0$ and passes
through exactly the matrices of Fig.~\ref{fig:knuth16rows}.

\begin{figure}[ht]
\centering
\begin{tikzpicture}[x=1.25cm,y=.9cm]
\node[girl,free] (g7) at (0,0) {$g_7$};
\node[boy] (b7) at (1,0) {$b_7$};
\node[girl] (g8) at (2,0) {$g_8$};
\node[boy,free] (b8) at (3,.75) {$b_8$};
\node[boy] (b5) at (3,-.75) {$b_5$};
\node[girl] (g6) at (4,-.75) {$g_6$};
\node[boy,free] (b4) at (5,-.75) {$b_4$};
\begin{scope}[on background layer]
\draw[sap] (g7.center)--(b7.center)--(g8.center)--(b8.center);
\draw[draw=boydraw,line width=6.5pt,line cap=round,line join=round,opacity=.22]
   (g7.center)--(b7.center)--(g8.center)--(b5.center)--(g6.center)--(b4.center);
\end{scope}
\draw[pedge] (g7)--(b7);
\draw[couple] (b7)--(g8);
\draw[pedge] (g8)--(b8);
\draw[edge,line width=.75pt,draw=black!60] (g8)--(b5);
\draw[couple] (b5)--(g6);
\draw[edge,line width=.75pt,draw=black!60] (g6)--(b4);
\node[font=\small,anchor=west] at (3.35,.75) {$\pi_2$: level 2};
\node[font=\small,anchor=west] at (5.35,-.75) {$\pi_3$: level 3};
\end{tikzpicture}

\caption{Part of Knuth's 16-pair graph after~\Hs1. Both $\pi_2$ (level~2,
three edges) and $\pi_3$ (level~3, five edges) are augmenting paths.}
\label{fig:knuth16paths}
\end{figure}

\subsection{Three readings of the definition}
\label{sec:readings}

Exercise~15(c) calls $G$ a Tchoukaillon graph if it might cause
Algorithm~H ``to continually satisfy assumption (a)'' while following the
schedule. Since the algorithm is deterministic, ``might'' must refer to
the input order of Section~\ref{sec:input}. The example of
Fig.~\ref{fig:knuth16paths} shows that the matrices of answer~15(d)
provide something weaker than assumption~(a). We therefore name three
properties that a round with $f$ free girls at level~$L$ may have:
\begin{itemize}
\item[$\Up$:] every augmenting path has level~$L$ (the literal hypothesis
  of 15(a); $\Up$ for ``uniform'');
\item[$\Bp$:] there are $f$ vertex-disjoint SAPs, that is, $\nu=f$;
  equivalently, for some perfect matching all $f$ basic paths are SAPs
  (what the matrices show; $\Bp$ for ``full'');
\item[$\Op$:] no condition beyond the numbers.
\end{itemize}

\begin{definition}
An ordered graph with a perfect matching \emph{realizes the Tchoukaillon
schedule} if~\Hs1 leaves $\floor{s/2}$ free girls and every later round
has the level~$L$ and the number of SAPs $a_L$ prescribed by
\eqref{eq:schedule}. A realization is \emph{literal} if every round after
\Hs1 has property~$\Up$, and \emph{basic-row} if every such round has
property~$\Bp$. We call a realization with no extra condition
\emph{operational}.
\end{definition}

Property $\Up$ implies $\Bp$, since the basic paths are augmenting paths.
Property $\Bp$ is exactly what Lemma~\ref{lem:15ab}(b) needs to force
$\ceil{f/(L+1)}$ SAPs. Neither implication reverses: Knuth's graph is a
basic-row realization that is not literal, and Section~\ref{sec:no-cover}
gives an operational realization (an 8-pair tree) in which a round has
two free girls but only one SAP, so it is not basic-row.

Which reading did Knuth intend? The literal one is what 15(a) and 15(c)
say. The matrices of 15(d) support the basic-row reading, and the
motivation---making Algorithm~H slow---supports the operational one;
the abstract of Knuth's 2022 lecture on Tchoukaillon numbers
\cite{KnuthRutgers} poses the matching question in terms of running
time. These notes do not choose. They answer the literal question
completely (Section~\ref{sec:literal}), and report partial results on the
other two (Section~\ref{sec:operational}).

Finally there is the question raised in the answer to exercise~16:
whether some graph forces \emph{more} rounds than a Tchoukaillon graph
would, ``by violating assumption~(a) and/or by finding more SAPs.'' Here
no condition is imposed on the run at all. Let
\[
  R(s)=\max\ \{\,\text{number of rounds of Algorithm~H}\,\},
\]
the maximum over all ordered bipartite graphs whose maximum matching has
size~$s$. (Neither a perfect matching nor any particular initial
matching is required.) Sections~\ref{sec:upper}--\ref{sec:comparison}
determine $R(s)$ exactly. Table~\ref{tab:readings} summarizes the
four questions and what these notes establish about them.

\begin{table}[ht]
\centering\small
\begin{tabular}{>{\raggedright\arraybackslash}p{.2\textwidth}
  >{\raggedright\arraybackslash}p{.34\textwidth}
  >{\raggedright\arraybackslash}p{.36\textwidth}}
\toprule
Question & Requirement & Answer here\\
\midrule
Literal & schedule, and $\Up$ in every round & for $s\ge1$, exists
  exactly when $s\in\{1,2,3,4,5,8,9\}$ (Corollary~L)\\[3pt]
Basic-row & schedule, and $\Bp$ in every round & explicit inputs for all
  $s\le57$ (Appendix~\ref{app:catalog}), including Knuth's graph; open in
  general\\[3pt]
Operational & schedule only & the same inputs; open in general; necessary
  conditions in Section~\ref{sec:operational}\\[3pt]
Most rounds & nothing & $R(s)=\floor{\sqrt{4s+1}}-1$ (Theorem~R);
  always $R(s)\ge T(s)$, with strict inequality first at $s=9$ and for
  all large $s$ (Section~\ref{sec:comparison})\\
\bottomrule
\end{tabular}
\caption{Four questions suggested by exercises~15 and 16.}
\label{tab:readings}
\end{table}
